% !TEX program = pdflatex
% Overleaf-compatible. Self-contained: only stock packages.
% FULL DRAFT with the three-environment experiments.
\documentclass[11pt]{article}

\usepackage[margin=1in]{geometry}
\usepackage{graphicx}
\usepackage{booktabs}
\usepackage{amsmath}
\usepackage{amssymb}
\usepackage{xcolor}
\usepackage{caption}
\usepackage{subcaption}
\usepackage{enumitem}
\usepackage{float}
\usepackage[round,authoryear]{natbib}
\usepackage[colorlinks=true, allcolors=blue!55!black]{hyperref}
\usepackage[section]{placeins}   % flush pending floats at each \section

% relax float placement so figures sit near their text instead of piling up
\renewcommand{\topfraction}{0.92}
\renewcommand{\bottomfraction}{0.85}
\renewcommand{\textfraction}{0.06}
\renewcommand{\floatpagefraction}{0.72}
\setcounter{topnumber}{3}
\setcounter{bottomnumber}{2}
\setcounter{totalnumber}{4}

\captionsetup{font=small, labelfont=bf}
\newcommand{\code}[1]{\texttt{#1}}
\newcommand{\figdir}{figures/belief_manifolds}
\newcommand{\claim}[2]{\smallskip\noindent\textbf{Claim #1.}\ \emph{#2}\smallskip}

\title{\textbf{Beliefs as Steerable Attractors:}\\[2pt]
\large The Geometry and Causal Structure of Evidence Integration\\
in Model-Free Reinforcement-Learning Agents}
\author{}
\date{}

\begin{document}
\maketitle
\vspace{-2.2em}

% =====================================================================
\begin{abstract}\noindent
Recurrent reinforcement-learning (RL) agents acting under partial observation
must maintain a \emph{belief} about the unobserved state of their environment.
It is by now routine to \emph{decode} such beliefs from an agent's hidden state,
but decodability alone says little about what the belief \emph{is} --- a static
register that stores a value, or a dynamical process --- or whether it
\emph{causes} the agent's behavior. We study belief in model-free RL agents as a
dynamical, causal object across three partially observed decision tasks that
differ only in the \emph{structure of the evidence} they present: a categorical
cue, a stream of accumulating left/right evidence, and noisy naturalistic
terrain. We find, using one analysis pipeline, that (i) the agent's belief
occupies a low-dimensional manifold in recurrent-state space whose geometry is
dictated by the task's evidence model --- a set of discrete attractors for
categorical evidence, and a one-dimensional line attractor for evidence
accumulation (net evidence linearly decodable at $R^2=1.00$); (ii) this belief
is an integrating attractor, not a register: it forms gradually, persists across
gaps in observation with no external cue, and revises in a lawful direction when
evidence is forced off-policy; and (iii) behavior is a graded, monotone function
of position on the belief manifold --- clamping the hidden state at points along
the manifold traces a smooth psychometric curve that gives \emph{full} control of
the agent's choice --- yet the \emph{same} intervention is behaviorally inert in a
control task where the belief is represented and even persistently implantable
but not decision-relevant. This double dissociation shows that causal
load-bearing is a property of the task's decision structure, not of
representational quality. Along the way we show that the choice of belief metric
matters: a categorical (Hellinger) embedding is faithful for unordered latents
but folds an ordered scalar latent into a misleading horseshoe, which a
magnitude-respecting metric repairs. Our accumulation task instantiates, in a
reward-trained agent with full observability and exact causal control, the
line-attractor-and-perturbation paradigm long studied in decision neuroscience.
\end{abstract}

% =====================================================================
\section{Introduction}

An agent acting under partial observation cannot see the state that determines
the consequences of its actions. To act well it must maintain a \emph{belief}: a
summary of its history sufficient for choosing actions \citep{kaelbling1998planning}.
In deep RL this belief is not designed but learned --- typically it lives,
implicitly, in the hidden state of a recurrent policy \citep{lambrechts2022recurrent,
mikulik2020meta}. A growing body of interpretability work shows that such beliefs
can be \emph{read out}: linear probes recover the latent variables an agent should
track, and recent work even recovers the predicted \emph{geometry} of the optimal
belief in sequence models \citep{shai2024transformers, li2023othello}. Decoding a
variable, however, leaves the scientifically central questions untouched. It does
not tell us what \emph{kind} of object the belief is --- a value latched into a
register, or a process unfolding in the agent's dynamics --- and, notoriously, it
does not tell us whether the decoded variable actually \emph{causes} the agent's
behavior or is merely correlated with it \citep{ravichander2021probing}.

We take these two questions --- \emph{what form?} and \emph{does it cause?} --- as
the subject of the paper, and study them in model-free RL agents, where belief is
a by-product of reward maximization rather than an explicitly trained objective.
Our stance is that a belief is best understood as a \emph{dynamical, causal}
object, and that the way to characterize it is to (i) look at the \emph{geometry}
traced out by the agent's recurrent states, (ii) analyze the \emph{dynamics} that
move states through that geometry, and (iii) \emph{intervene} on those states and
measure the effect on behavior. This viewpoint is standard in the
systems-neuroscience study of decision circuits, where trained recurrent networks
and neural populations alike are reverse-engineered into low-dimensional dynamical
pictures --- fixed points, line attractors, and their perturbations
\citep{mante2013context, sussillo2013opening, wong2006recurrent}. We bring it to
bear on model-free RL agents, where we can do what neuroscience cannot: observe
the belief with ground truth and steer it with surgical precision.

The key design choice is to hold the analysis fixed while varying the
\emph{structure of the evidence} a task provides (Figure~\ref{fig:tasks}). In a
\emph{categorical-cue T-maze}, a single symbolic cue seen early must be carried to
a decision made later --- the belief is a discrete class. In an \emph{accumulation
T-maze} --- the perceptual decision-making paradigm of \citet{brunton2013rats} and
\citet{pinto2018towers}, rendered as navigation --- a stream of left/right
evidence must be integrated and the majority side chosen; the belief is a
\emph{scalar} accumulator. In a \emph{naturalistic navigation} task, the belief
(which terrain regime the agent is in) must be inferred from noisy, distributed
observations. Because the tasks differ only in evidence structure, differences in
the recovered belief object attach to that structure, and commonalities attest to
the method's generality.

\paragraph{Contributions.} We make three claims about belief in model-free RL
agents, each supported by the same pair of tools across all three tasks.

\claim{1 (Geometry follows evidence structure)}{The belief occupies a
low-dimensional manifold in recurrent-state space whose geometry is dictated by
the task's evidence model: discrete attractor points for categorical evidence and
a one-dimensional line attractor for scalar accumulation. One task-agnostic
construction recovers the correct object in each case.}

\claim{2 (Belief is an integrating attractor, not a register)}{The belief forms
gradually as evidence arrives, persists across an observation gap with no external
cue to sustain it, and revises in a lawful direction when fresh evidence is forced
off-policy; the recurrent dynamics show marginally stable slow manifolds
consistent with integration rather than a latched value.}

\claim{3 (Behavior is a graded causal function of manifold position, and this is
task-contingent)}{Clamping the hidden state at fractional positions along the
belief manifold produces a smooth psychometric curve over choice with full control
at the endpoints. Yet in a control task where the same latent is represented ---
and can even be persistently implanted --- but is not decision-relevant, the
identical intervention is behaviorally inert: a double dissociation showing that
causal load-bearing is a property of the task's decision structure, not of
decodability.}

\paragraph{Why it matters.} For interpretability, Claim~3 upgrades ``we can decode
$X$'' to ``$X$ is a causal, continuously manipulable variable,'' and gives a
graded test --- the manifold-clamp psychometric --- for whether a probed feature is
load-bearing, sharper than an all-or-nothing patch. For the science of learned
agents, viewing belief as attractor dynamics explains why memory persists across
observation gaps and predicts how it fails. And because our accumulation task
reproduces, in a reward-trained agent, the line-attractor mechanism and
perturbation logic foundational in decision neuroscience \citep{mante2013context,
hanks2015distinct}, it offers an in-silico model organism where those experiments
can be run with full observability and exact causal control.

% =====================================================================
\section{Related work}

\paragraph{Belief representations in RL and sequence models.} That recurrent POMDP
agents encode belief-like sufficient statistics is established theoretically and
empirically \citep{kaelbling1998planning, lambrechts2022recurrent}, and meta-RL
agents can approximate Bayes-optimal updates \citep{mikulik2020meta}. Recent
interpretability decodes such latents and even recovers the \emph{geometry} of the
optimal belief predicted by computational mechanics in transformers
\citep{shai2024transformers}, echoing emergent world-state representations in
game-playing models \citep{li2023othello}. We depart in two ways: we study
\emph{model-free} agents where belief is incidental to reward, and we go past
decoding to \emph{causal, graded} intervention and to \emph{dynamical}
characterization.

\paragraph{Dynamics of decision-making in trained RNNs and brains.} A long
tradition reverse-engineers recurrent computation into low-dimensional dynamics
--- fixed points and line attractors underlying integration and choice
\citep{sussillo2013opening, mante2013context, wong2006recurrent}, with evidence
accumulation as the canonical case \citep{gold2007neural, brunton2013rats}. Causal
perturbations of the accumulator in cortex sharpened this picture
\citep{hanks2015distinct}, and belief-state accounts connect it to reward learning
\citep{starkweather2017dopamine}. Our accumulation environment is a deliberate RL
analogue of the accumulating-towers task \citep{pinto2018towers}; our contribution
is to show the same dynamical object arising in a \emph{model-free, reward-trained}
agent and to steer it along its manifold.

\paragraph{Activation-level interpretation and steering.} Causal mediation via
activation patching and representation-level steering vectors are common tools for
locating and manipulating features \citep{meng2022locating, turner2023activation,
zou2023representation}, and probing studies caution that decodability need not
imply causal use \citep{ravichander2021probing}. Rather than restate that caution,
we operationalize its resolution: a manifold-clamp psychometric that measures
\emph{how much} of the behavior a belief controls, and a matched-representation
control task that isolates \emph{when} a represented belief is load-bearing.

% =====================================================================
\section{Setup and definitions}
\label{sec:setup}

\paragraph{Partially observed decision tasks.} Each task is a POMDP in which a
hidden latent $z$ --- the categorical cue, the sign/magnitude of accumulated
evidence, or the terrain regime --- determines which terminal choice is rewarded.
The agent receives egocentric symbolic observations $o_t$ and emits actions
$a_t$; a recurrent policy $\pi_\theta(a_t \mid h_t)$ carries a hidden state
$h_t = f_\theta(h_{t-1}, o_t)$ (a GRU). All three tasks share an observation and
action interface, so a single architecture and analysis pipeline apply unchanged.
We train model-free recurrent agents with PPO \citep{schulman2017proximal} to high
success and freeze them for analysis; the accumulation agent reaches $100\%$
success (optimal-length trajectories), and the cue agents reach $100\%$.

\paragraph{Belief.} Operationally, the belief at time $t$ is a distribution
$b_t = g(h_t)$ over the task-relevant latent, where $g$ is a \emph{linear} read-out
(a multinomial logistic probe) fit on held-out episodes from the frozen hidden
states. Linearity keeps the belief an affine function of $h_t$, so geometry in
belief space corresponds to geometry in hidden space and interventions can be
expressed in either. We stress that $b_t$ is a \emph{measurement} of the agent's
state, not a claim that the agent computes it explicitly; the causal experiments
(\S\ref{sec:e3}) license statements about what the belief \emph{does}.

% =====================================================================
\section{Method}
\label{sec:method}
One triplet of tools, applied identically to every task.

\paragraph{(1) Belief-space geometry.} From many frozen rollouts we collect hidden
states $\{h_t\}$ and probe beliefs $\{b_t\}$. To visualize the shape of the set of
beliefs we embed them with a divergence-preserving projection: map each belief
through the Hellinger embedding $b \mapsto \sqrt{b}$ and take leading principal
components (equivalently, information-geometric MDS on the Bhattacharyya
divergence). Confident beliefs land at simplex vertices and ambiguous beliefs in
the interior, so the number of distinguishable states and how the agent moves
between them become directly visible and comparable across tasks. For
\emph{ordered} latents we additionally use a magnitude-respecting embedding
(\S\ref{sec:metric}).

\paragraph{(2) Attractor analysis.} To ask whether the belief is a dynamical
attractor we analyze the frozen recurrent map $h \mapsto f_\theta(h,o)$ under a
fixed post-evidence observation, locating slow points by minimizing the update
magnitude $\lVert f_\theta(h,o) - h\rVert^2$ and reading local stability from the
Jacobian spectrum. Marginally stable directions (eigenvalues near the unit circle)
signal integration/memory --- a continuum of near-fixed states along which evidence
pushes the belief and where it rests.

\paragraph{(3) Manifold-clamp steering.} To test causal load-bearing in a graded
way we intervene on the hidden state. Given the belief manifold's local axis
$\hat u(x)$ at task-phase $x$ (the direction between class-conditional means, made
position-local because the code can rotate along a trajectory), we clamp the
component of $h_t$ along $\hat u$ to a target $\rho_f$ interpolating between the
endpoints at fraction $f\in[0,1]$:
\[
  h_t \;\leftarrow\; h_t + \big(\rho_f - \langle h_t, \hat u\rangle\big)\,\hat u,
\]
sustained through the decision phase. Sweeping $f$ and reading out the terminal
choice yields a \emph{psychometric curve}: behavior as a function of implanted
belief. A smooth, monotone curve with full endpoint control means the belief is
continuously steerable and causally load-bearing; a flat curve means it is not.

% =====================================================================
\section{Tasks}
\label{sec:tasks}

\begin{figure}[tbp]
  \centering
  \begin{subfigure}{0.62\linewidth}
    \includegraphics[width=\linewidth]{\figdir/env_towers_maps.png}
    \caption{accumulation T-maze: a selection of maps}
  \end{subfigure}\hfill
  \begin{subfigure}{0.36\linewidth}
    \includegraphics[width=\linewidth]{\figdir/env_towers_filmstrip.png}
    \caption{egocentric evidence stream}
  \end{subfigure}
  \caption{\textbf{The accumulation T-maze (\code{towers}).} The agent runs east;
  a tower on the top wall is left evidence, on the bottom wall right evidence.
  After the evidence ends and leaves view, a T-junction requires choosing the
  majority side (white ring). (b) The egocentric $5{\times}5$ view: towers enter
  and leave as discrete evidence pulses while the running accumulator climbs. The
  three environments share this observation/action interface.}
  \label{fig:tasks}
\end{figure}

The three environments span the axis of evidence structure while sharing
interface, architecture, and analysis.

\paragraph{Categorical cue (\code{memory}; discrete belief).} A T-maze in which a
single cue, visible only briefly near the start, specifies the rewarded
branch/door; the belief is one of $K$ classes. We use $K\in\{2,3,4\}$ cue sets. A
cue carries two attributes (a shape indicating the branch and a colour indicating
the door); the $2$-cue set entangles them (two cues) while the $4$-cue set
factorizes them (all shape$\times$colour combinations), which lets us test whether
the agent's belief factorizes. The cue leaves view long before the choice, so the
task exercises memory (persistence) most starkly.

\paragraph{Evidence accumulation (\code{towers}; scalar belief).} The
neuroscience-canonical case as navigation (Figure~\ref{fig:tasks}): a stream of
left/right ``towers'' must be integrated and the \emph{majority} side chosen after
the evidence leaves view. The belief is a scalar accumulator (net right$-$left).

\paragraph{Naturalistic navigation (\code{bridge\_tunnel}; noisy belief).} A
procedurally generated terrain-crossing task in which the belief --- the terrain
regime (rocky / lakes / balanced) that dictates the correct strategy --- must be
inferred from noisy, spatially distributed observations. It supplies the crucial
\emph{control}: two variants represent the same belief but differ in whether it is
decision-relevant. In \code{fork\_wall}, terrain determines which of two doors is
correct (belief is load-bearing); in \code{commit}, the agent commits to a
crossing strategy where the belief need not gate the choice.

% =====================================================================
\section{Experiment 1: belief geometry follows evidence structure}
\label{sec:e1}

\begin{figure}[tbp]
  \centering
  \includegraphics[width=0.85\linewidth]{\figdir/geom_memory_vertices.png}\\[3pt]
  \includegraphics[width=\linewidth]{\figdir/geom_towers_line.png}
  \caption{\textbf{Belief geometry is dictated by evidence structure.}
  \emph{Top:} the \code{memory} categorical belief embeds as \emph{discrete
  vertices} of a probability simplex (here the $3$/$4$-cue sets); confident beliefs
  are mutually equidistant and the number of occupied vertices equals the number
  of cue classes. \emph{Bottom:} the \code{towers} scalar belief is a
  \emph{one-dimensional line attractor}: net evidence is linearly decoded from the
  GRU state at $R^2=1.00$ (C); post-evidence states form a single axis ordered by
  net evidence (A, correlation $0.92$ with the leading PC); the belief forms along
  this axis over time (B); and the belief read-out is the sharp step of a
  near-perfect integrator (D).}
  \label{fig:geometry}
\end{figure}

We construct the belief-space embedding (\S\ref{sec:method}) for each frozen agent
(Figure~\ref{fig:geometry}). The geometries differ exactly as the evidence
structure predicts, from one construction.

\paragraph{Categorical $\Rightarrow$ discrete attractors.} For \code{memory}, the
belief embeds as isolated vertices of a probability simplex, one per cue class:
$K$ classes occupy $K$ vertices spanning $K{-}1$ dimensions (with a small
residual). Confident beliefs are mutually equidistant, so the class \emph{identity}
--- not a magnitude --- is what the geometry encodes.

\paragraph{Scalar $\Rightarrow$ line attractor.} For \code{towers}, the belief is a
one-dimensional line. Net accumulated evidence is decoded from the GRU state by a
single linear map at $R^2=1.00$; post-evidence states arrange along one dominant
axis ordered by net evidence (leading principal component correlation $0.92$); and
early states start near the center and sweep outward as evidence arrives ---
integration, not a latched value. The attractor analysis (\S\ref{sec:e2}) confirms
the axis is a slow manifold.

\paragraph{Naturalistic $\Rightarrow$ a graded, noisier manifold.} For
\code{bridge\_tunnel}, the same construction yields a category belief that forms
gradually from distributed terrain evidence, with the interior of the simplex
genuinely occupied during inference (Figure~\ref{fig:bt_geometry}). The belief is
weaker in a plain agent (held-out category decode $\approx0.48$) and sharpened by
an auxiliary belief head ($\approx0.76$), but present in both.

\begin{figure}[tbp]
  \centering
  \includegraphics[width=0.92\linewidth]{\figdir/geom_bt.png}
  \caption{\textbf{Naturalistic belief (\code{bridge\_tunnel}).} The map-category
  belief (rocky/lakes/balanced) inferred from noisy terrain, embedded with the same
  Hellinger construction and evaluated on held-out maps; belief forms gradually
  over the episode (right) rather than snapping to a vertex.}
  \label{fig:bt_geometry}
\end{figure}

\subsection{The belief metric must match the latent's structure}
\label{sec:metric}

The Hellinger/Bhattacharyya embedding treats belief classes as \emph{unordered}.
This is faithful for the categorical cue --- there is no magnitude to represent, and
the equidistant simplex vertices are the honest picture. For the \emph{ordered}
scalar of \code{towers}, however, it is actively misleading: the Hellinger distance
\emph{saturates} at $\sqrt2$ as soon as two beliefs stop overlapping, so the
maximally opposite beliefs (net $+4$ vs $-4$) sit at the same ceiling distance as
moderately different ones, and a linear projection of this bounded, ordered metric
folds into a \emph{horseshoe} in which the extremes curl together
(Figure~\ref{fig:metric}A; projected distance $0.08$ for the opposite extremes
versus $0.15$ for adjacent classes --- an inversion of the truth). A casual reader
would misread the folded extremes as similar beliefs.

Two remedies recover a faithful picture (Figure~\ref{fig:metric}B,C): a
magnitude-respecting distribution metric (Wasserstein-1, the $L_1$ distance of CDFs
on the ordered support), whose classical MDS lays the beliefs on a straight line
with the extremes maximally far ($d(+4,-4)=7.76$ vs.\ $d(+4,+3)=0.90$, all variance
on one axis); or simply the scalar read-out itself, $\mathbb{E}[\text{net}\mid b]$
against belief confidence, which cannot fold. The methodological point generalizes:
\emph{match the belief metric to the latent's structure --- categorical
$\rightarrow$ Hellinger; ordinal/scalar $\rightarrow$ Wasserstein or the raw scalar
axis.}

\begin{figure}[tbp]
  \centering
  \includegraphics[width=\linewidth]{\figdir/metric_faithful.png}
  \caption{\textbf{The metric encodes an assumption about the latent.} (A)
  Hellinger$+$PCA of the ordered scalar belief folds into a horseshoe: the opposite
  extremes $\pm4$ project close together (misleading). (B) Wasserstein$+$MDS lays
  the same beliefs on a straight ordered line with the extremes maximally far. (C)
  The scalar read-out $\mathbb{E}[\text{net}\mid b]$ vs.\ confidence is faithful by
  construction.}
  \label{fig:metric}
\end{figure}

% =====================================================================
\section{Experiment 2: beliefs are integrating attractors, not registers}
\label{sec:e2}

\paragraph{Formation and persistence.} In all tasks the belief \emph{forms}
gradually --- it moves along the manifold from ambiguity toward a
vertex/endpoint as evidence arrives (Figure~\ref{fig:geometry}B for \code{towers};
the cue tasks snap to a vertex once the cue is seen and the naturalistic task ramps
over many steps). Crucially the belief then \emph{persists} across an observation
gap with nothing sustaining it: in \code{towers} the accumulator is held stably
through the run-out corridor after all towers leave view, and in \code{memory} the
cue class is held from the cue room to the distant door. The attractor analysis of
the frozen recurrent map finds no exact fixed points but hundreds of slow points
with leading Jacobian eigenvalue magnitude $\approx 1.000$ --- approximately
marginally stable slow manifolds (one per belief fiber), the dynamical signature of
integration/memory rather than a latched register.

\paragraph{Lawful revision under forced evidence.} We test integration causally by
forcing the agent, off-policy, through evidence that contradicts its current
belief and reading out the resulting belief (Figure~\ref{fig:mem_evidence}). In
\code{memory}, forcing the $2$-cue agent down the wrong corridor \emph{collapses its
belief to the opposite attractor} --- it comes to ``believe'' the other cue and its
downstream door choice flips accordingly; the $4$-cue agent instead \emph{revises
only the branch/direction factor} while preserving the colour factor, and the
$3$-cue agent revises in the direction consistent with its training likelihoods.
The belief is thus not a frozen record of the cue but a state that integrates new
evidence.

\begin{figure}[tbp]
  \centering
  \includegraphics[width=0.9\linewidth]{\figdir/mem_evidence.png}
  \caption{\textbf{Beliefs integrate forced evidence (\code{memory}, $2$-cue).}
  Forcing the agent through the wrong corridor drives its decoded belief from one
  attractor to the other (aligned to the moment the marker is opened), and its
  downstream choice follows --- the belief is updated by evidence, not merely stored.}
  \label{fig:mem_evidence}
\end{figure}

% =====================================================================
\section{Experiment 3: behavior is a graded causal function of manifold position}
\label{sec:e3}

The central experiment clamps the hidden state at fractions $f$ along the belief
manifold at the decision phase and reads out the choice (\S\ref{sec:method}).

\paragraph{\code{towers}: full graded control, evidence-independent.} Clamping the
accumulator belief traces a clean psychometric (Figure~\ref{fig:psy_towers}B):
$P(\text{right door})$ moves $0.00 \!\to\! 0.03\,(f{=}0.5) \!\to\! 0.97\,(f{=}0.62)
\!\to\! 1.00$, with \emph{zero} timeouts --- full control at both endpoints. The
transition is steep (the sharp sigmoid of a near-perfect integrator). Splitting by
the map's true majority (Figure~\ref{fig:psy_towers}C), true-left and true-right
maps follow the \emph{same} curve: the implanted belief \emph{overrides} the
accumulated evidence, and the decision is a single path-independent function of the
belief coordinate. The Hellinger belief arc (A) is shown for continuity with the
other tasks; per \S\ref{sec:metric} it is the saturated image of the line attractor.

\begin{figure}[tbp]
  \centering
  \includegraphics[width=\linewidth]{\figdir/psy_towers.png}
  \caption{\textbf{Manifold-clamp psychometric (\code{towers}).} (A) Hellinger
  belief space (ordered arc). (B) Clamping the GRU belief along the manifold gives
  full graded control of the door with no timeouts. (C) True-left and true-right
  maps follow one curve: the clamp overrides the evidence (path-independent
  boundary).}
  \label{fig:psy_towers}
\end{figure}

\paragraph{\code{memory}: the psychometric respects the belief's factorization.}
The same clamp along the cue-belief manifold (Figure~\ref{fig:psy_memory}) moves the
probe belief smoothly and the behavior as a steep sigmoid. In the entangled $2$-cue
agent the branch and door choices are \emph{yoked} --- both flip together at the same
$f$, because a single cue code carries both attributes. In the factorized $4$-cue
agent, clamping the shape/direction axis flips the branch smoothly while the door
stays \emph{invariant} at $P=1.0$ across the entire edge, because the colour factor
--- and thus the door --- is left untouched. The causal geometry of the intervention
mirrors the representational geometry of the belief.

\begin{figure}[tbp]
  \centering
  \includegraphics[width=0.95\linewidth]{\figdir/mem_psychometric.png}
  \caption{\textbf{Manifold-clamp psychometric (\code{memory}).} Probe belief (grey)
  moves smoothly with clamp position; behavior is a steep sigmoid. \emph{Left,}
  $4$-cue: clamping the shape axis flips the branch (purple) while the door (green)
  stays at $1.0$ --- the intervention respects the belief's factorization.
  \emph{Right,} $2$-cue: branch and door are yoked and flip together.}
  \label{fig:psy_memory}
\end{figure}

\paragraph{\code{bridge\_tunnel} \code{fork\_wall}: graded control in a noisy task.}
The same clamp on the naturalistic category belief gives a graded, monotone
psychometric over the door choice (Figure~\ref{fig:psy_bt}), with full control at
$f{=}1$ in both sweep directions. Sweeping from lakes toward rocky the choice
crosses $50/50$ at $f\approx0.62$; sweeping from rocky toward lakes it crosses at
$f\approx0.38$ --- \emph{the same physical point} ($\approx\!62\%$ toward rocky),
a single path-independent boundary. The transition is broader than in
\code{towers}/\code{memory}, consistent with ambient terrain evidence partially
contradicting the implant at intermediate doses. Belief is load-bearing here.

\begin{figure}[tbp]
  \centering
  \includegraphics[width=0.95\linewidth]{\figdir/psy_bt_forkwall.png}
  \caption{\textbf{Manifold-clamp psychometric (\code{bridge\_tunnel}
  \code{fork\_wall}).} Door choice varies smoothly with clamp position and is fully
  controlled at $f{=}1$ in both directions; the $50/50$ boundary is the same physical
  point from either sweep ($\approx\!62\%$ toward rocky) --- path-independent.}
  \label{fig:psy_bt}
\end{figure}

\paragraph{\code{bridge\_tunnel} \code{commit}: the belief is epiphenomenal ---
a double dissociation.} In the \code{commit} variant the category belief is present
and decodable (auxiliary agent held-out $\approx0.76$) and can even be
\emph{persistently implanted}: a full belief swap during an early window survives
long after release, the decoded belief staying on the wrong category for $100^{+}$
steps against contradicting terrain (Figure~\ref{fig:psy_commit}). Yet the agent's
behavior is \emph{unchanged} --- its commitment matrix and success rate are identical
to baseline. The identical intervention that fully controls choice in
\code{fork\_wall} is behaviorally inert in \code{commit}. Representational quality,
decodability, and even persistence do not make a belief causal; the task's decision
structure does.

\begin{figure}[tbp]
  \centering
  \includegraphics[width=0.95\linewidth]{\figdir/psy_bt_commit.png}
  \caption{\textbf{Persistent but epiphenomenal belief (\code{bridge\_tunnel}
  \code{commit}).} A full belief swap (dotted) during the clamp window sticks after
  release --- the decoded belief remains on the wrong category for $100^{+}$ steps ---
  while a half-belief (dashed) re-corrects. Despite the persistent implant, the
  agent's commitment behavior is unchanged from baseline (text): the belief is
  represented but not decision-relevant.}
  \label{fig:psy_commit}
\end{figure}

\paragraph{Summary of the causal picture.} Table~\ref{tab:summary} collects the
three environments. The same manifold-clamp intervention gives three transfer
functions on one axis --- steep-full, factorized, graded-full, and flat --- selected
by the task's decision structure, not by how well the belief can be read out.

\begin{table}[tbp]
\centering\small
\begin{tabular}{@{}lllll@{}}
\toprule
Environment & Latent & Belief geometry & Clamp $\rightarrow$ behavior & Belief causal? \\
\midrule
\code{memory} ($2$-cue)  & categorical (entangled)  & simplex vertices & steep; branch+door yoked & yes \\
\code{memory} ($4$-cue)  & categorical (factorized) & simplex vertices & steep; branch flips, door fixed & yes (factored) \\
\code{towers}            & scalar accumulator       & line attractor   & steep, full, evidence-overriding & yes \\
\code{bridge} fork\_wall & category (noisy)         & graded simplex   & graded, full, path-independent & yes \\
\code{bridge} commit     & category (noisy)         & graded simplex   & flat (implant persists, inert) & no \\
\bottomrule
\end{tabular}
\caption{One intervention, four transfer functions, selected by decision
structure. Decodability is high in all rows; causal load-bearing is not.}
\label{tab:summary}
\end{table}

% =====================================================================
\section{Discussion}

Across three tasks that differ only in evidence structure, a model-free RL agent
represents its task-relevant belief as a \emph{dynamical attractor manifold} whose
geometry mirrors the task's evidence model, integrates evidence over time, persists
without external support, and --- where the task makes it decision-relevant --- is a
graded, continuously steerable cause of behavior. The manifold-clamp psychometric
turns ``is this feature causal?'' into a quantitative dose--response curve, and the
\code{commit} control shows why that matters: decodability, representational
sharpness, and even implant persistence can all be present in a belief that does
nothing. We also found that the \emph{metric} used to visualize belief is a modeling
choice with teeth --- a categorical embedding folds an ordered latent into a
misleading horseshoe --- and gave a simple rule for matching metric to latent.

\paragraph{Scope and limitations.} The agents are small recurrent policies on
gridworld-scale POMDPs, chosen for the ground truth and causal access they afford;
claims about ``model-free RL agents'' should be read as claims about this regime,
not about frontier-scale agents. The belief is a linear measurement of the hidden
state; non-linear structure it misses is a caveat we probe but cannot fully exclude.
``Integration'' and the neuroscience analogy are about consistency with the task's
evidence statistics up to the training distribution, not optimal Bayesian inference.
Finally, the manifold-clamp axis is estimated from class-conditional means; we make
it position-local to handle rotation of the code along a trajectory, but a
mis-estimated axis would under-state causal control (a conservative bias).

% =====================================================================
\begin{thebibliography}{99}\small
\bibitem[Brunton et al.(2013)]{brunton2013rats} B.~Brunton, M.~Botvinick, and
C.~Brody. Rats and humans can optimally accumulate evidence for decision-making.
\emph{Science}, 2013.
\bibitem[Gold \& Shadlen(2007)]{gold2007neural} J.~Gold and M.~Shadlen. The neural
basis of decision making. \emph{Annual Review of Neuroscience}, 2007.
\bibitem[Hanks et al.(2015)]{hanks2015distinct} T.~Hanks, C.~Kopec, B.~Brunton,
C.~Duan, J.~Erlich, and C.~Brody. Distinct relationships of parietal and prefrontal
cortices to evidence accumulation. \emph{Nature}, 2015.
\bibitem[Kaelbling et al.(1998)]{kaelbling1998planning} L.~Kaelbling, M.~Littman,
and A.~Cassandra. Planning and acting in partially observable stochastic domains.
\emph{Artificial Intelligence}, 1998.
\bibitem[Lambrechts et al.(2022)]{lambrechts2022recurrent} G.~Lambrechts,
A.~Bolland, and D.~Ernst. Recurrent networks, hidden states and beliefs in
partially observable environments. \emph{TMLR}, 2022.
\bibitem[Li et al.(2023)]{li2023othello} K.~Li, A.~K. Hopkins, D.~Bau, F.~Vi\'egas,
H.~Pfister, and M.~Wattenberg. Emergent world representations. \emph{ICLR}, 2023.
\bibitem[Mante et al.(2013)]{mante2013context} V.~Mante, D.~Sussillo, K.~Shenoy,
and W.~Newsome. Context-dependent computation by recurrent dynamics in prefrontal
cortex. \emph{Nature}, 2013.
\bibitem[Meng et al.(2022)]{meng2022locating} K.~Meng, D.~Bau, A.~Andonian, and
Y.~Belinkov. Locating and editing factual associations in GPT. \emph{NeurIPS}, 2022.
\bibitem[Mikulik et al.(2020)]{mikulik2020meta} V.~Mikulik, G.~Del\'etang,
T.~McGrath, et al. Meta-trained agents implement Bayes-optimal agents.
\emph{NeurIPS}, 2020.
\bibitem[Pinto et al.(2018)]{pinto2018towers} L.~Pinto, S.~A. Koay, et al. An
accumulation-of-evidence task using visual pulses for mice navigating in virtual
reality. \emph{Frontiers in Behavioral Neuroscience}, 2018.
\bibitem[Ravichander et al.(2021)]{ravichander2021probing} A.~Ravichander,
Y.~Belinkov, and E.~Hovy. Probing the probing paradigm. \emph{EACL}, 2021.
\bibitem[Schulman et al.(2017)]{schulman2017proximal} J.~Schulman, F.~Wolski,
P.~Dhariwal, A.~Radford, and O.~Klimov. Proximal policy optimization algorithms.
\emph{arXiv:1707.06347}, 2017.
\bibitem[Shai et al.(2024)]{shai2024transformers} A.~Shai, S.~Marzen, P.~Riechers,
et al. Transformers represent belief state geometry in their residual stream.
\emph{arXiv}, 2024.
\bibitem[Starkweather et al.(2017)]{starkweather2017dopamine} C.~Starkweather,
B.~Babayan, N.~Uchida, and S.~Gershman. Dopamine reward prediction errors reflect
hidden-state inference across time. \emph{Nature Neuroscience}, 2017.
\bibitem[Sussillo \& Barak(2013)]{sussillo2013opening} D.~Sussillo and O.~Barak.
Opening the black box: low-dimensional dynamics in high-dimensional recurrent neural
networks. \emph{Neural Computation}, 2013.
\bibitem[Turner et al.(2023)]{turner2023activation} A.~Turner, L.~Thiergart, et al.
Activation addition: steering language models without optimization. \emph{arXiv},
2023.
\bibitem[Wong \& Wang(2006)]{wong2006recurrent} K.-F. Wong and X.-J. Wang. A
recurrent network mechanism of time integration in perceptual decisions.
\emph{Journal of Neuroscience}, 2006.
\bibitem[Zou et al.(2023)]{zou2023representation} A.~Zou, L.~Phan, S.~Chen, et al.
Representation engineering: a top-down approach to AI transparency. \emph{arXiv},
2023.
\end{thebibliography}

\end{document}
